\documentclass[
  doc,
  longtable,
  a4paper,
  nolmodern,
  notxfonts,
  notimes,
  colorlinks=true,linkcolor=blue,citecolor=blue,urlcolor=blue]{apa7}

\usepackage{amsmath}
\usepackage{amssymb}
\usepackage{setspace}
\geometry{inner=1in, outer=1in}
\fancyhfoffset[LE,RO]{0cm}



\RequirePackage{longtable}
\RequirePackage{threeparttablex}

\usepackage{framed}


\makeatletter
\renewcommand{\paragraph}{\@startsection{paragraph}{4}{\parindent}%
	{0\baselineskip \@plus 0.2ex \@minus 0.2ex}%
	{-.5em}%
	{\normalfont\normalsize\bfseries\typesectitle}}

\renewcommand{\subparagraph}[1]{\@startsection{subparagraph}{5}{0.5em}%
	{0\baselineskip \@plus 0.2ex \@minus 0.2ex}%
	{-\z@\relax}%
	{\normalfont\normalsize\bfseries\itshape\hspace{\parindent}{#1}\textit{\addperi}}{\relax}}
\makeatother




\usepackage{longtable, booktabs, multirow, multicol, colortbl, hhline, caption, array, float, xpatch}
\usepackage{subcaption}


\renewcommand\thesubfigure{\Alph{subfigure}}
\setcounter{topnumber}{2}
\setcounter{bottomnumber}{2}
\setcounter{totalnumber}{4}
\renewcommand{\topfraction}{0.85}
\renewcommand{\bottomfraction}{0.85}
\renewcommand{\textfraction}{0.15}
\renewcommand{\floatpagefraction}{0.7}

\usepackage{tcolorbox}
\tcbuselibrary{listings,theorems, breakable, skins}
\usepackage{fontawesome5}

\definecolor{quarto-callout-color}{HTML}{909090}
\definecolor{quarto-callout-note-color}{HTML}{0758E5}
\definecolor{quarto-callout-important-color}{HTML}{CC1914}
\definecolor{quarto-callout-warning-color}{HTML}{EB9113}
\definecolor{quarto-callout-tip-color}{HTML}{00A047}
\definecolor{quarto-callout-caution-color}{HTML}{FC5300}
\definecolor{quarto-callout-color-frame}{HTML}{ACACAC}
\definecolor{quarto-callout-note-color-frame}{HTML}{4582EC}
\definecolor{quarto-callout-important-color-frame}{HTML}{D9534F}
\definecolor{quarto-callout-warning-color-frame}{HTML}{F0AD4E}
\definecolor{quarto-callout-tip-color-frame}{HTML}{02B875}
\definecolor{quarto-callout-caution-color-frame}{HTML}{FD7E14}

%\newlength\Oldarrayrulewidth
%\newlength\Oldtabcolsep


\usepackage{hyperref}
\usepackage[protrusion=true]{microtype}



\providecommand{\tightlist}{%
  \setlength{\itemsep}{0pt}\setlength{\parskip}{0pt}}
\usepackage{longtable,booktabs,array}
\newcounter{none} % for unnumbered tables
\usepackage{calc} % for calculating minipage widths
% Correct order of tables after \paragraph or \subparagraph
\usepackage{etoolbox}
\makeatletter
\patchcmd\longtable{\par}{\if@noskipsec\mbox{}\fi\par}{}{}
\makeatother
% Allow footnotes in longtable head/foot
\IfFileExists{footnotehyper.sty}{\usepackage{footnotehyper}}{\usepackage{footnote}}
\makesavenoteenv{longtable}

\usepackage{graphicx}
\makeatletter
\newsavebox\pandoc@box
\newcommand*\pandocbounded[1]{% scales image to fit in text height/width
  \sbox\pandoc@box{#1}%
  \Gscale@div\@tempa{\textheight}{\dimexpr\ht\pandoc@box+\dp\pandoc@box\relax}%
  \Gscale@div\@tempb{\linewidth}{\wd\pandoc@box}%
  \ifdim\@tempb\p@<\@tempa\p@\let\@tempa\@tempb\fi% select the smaller of both
  \ifdim\@tempa\p@<\p@\scalebox{\@tempa}{\usebox\pandoc@box}%
  \else\usebox{\pandoc@box}%
  \fi%
}
% Set default figure placement to htbp
\def\fps@figure{htbp}
\makeatother







\usepackage{newtx}

\defaultfontfeatures{Scale=MatchLowercase}
\defaultfontfeatures[\rmfamily]{Ligatures=TeX,Scale=1}





\title{Importing Data in R: Exploring Various Packages and Their
Advantages}


\shorttitle{Importing Data in R: Exploring Various Packages and Their
Advantages}


\usepackage{etoolbox}


\course{Data Analysis for Decision-Making}
\professor{Prof.~Dr.~Stephan Huber}




\author{Ali Soleimani}



\affiliation{
{Fresenius University of Applied Science}}




\leftheader{Soleimani}

\date{May 8, 2026}





\authornote{ 

\par{       }
\par{Correspondence concerning this article should be addressed to Ali
Soleimani, Email: \href{mailto:soleimani.ali@stud.hs-fresenius.de}{soleimani.ali@stud.hs-fresenius.de}}
}

\makeatletter
\let\endoldlt\endlongtable
\def\endlongtable{
\hline
\endoldlt
}
\makeatother

\urlstyle{same}



\makeatletter
\@ifpackageloaded{caption}{}{\usepackage{caption}}
\AtBeginDocument{%
\ifdefined\contentsname
  \renewcommand*\contentsname{Table of contents}
\else
  \newcommand\contentsname{Table of contents}
\fi
\ifdefined\listfigurename
  \renewcommand*\listfigurename{List of Figures}
\else
  \newcommand\listfigurename{List of Figures}
\fi
\ifdefined\listtablename
  \renewcommand*\listtablename{List of Tables}
\else
  \newcommand\listtablename{List of Tables}
\fi
\ifdefined\figurename
  \renewcommand*\figurename{Figure}
\else
  \newcommand\figurename{Figure}
\fi
\ifdefined\tablename
  \renewcommand*\tablename{Table}
\else
  \newcommand\tablename{Table}
\fi
}
\@ifpackageloaded{float}{}{\usepackage{float}}
\floatstyle{ruled}
\@ifundefined{c@chapter}{\newfloat{codelisting}{h}{lop}}{\newfloat{codelisting}{h}{lop}[chapter]}
\floatname{codelisting}{Listing}
\newcommand*\listoflistings{\listof{codelisting}{List of Listings}}
\makeatother
\makeatletter
\makeatother
\makeatletter
\@ifpackageloaded{caption}{}{\usepackage{caption}}
\@ifpackageloaded{subcaption}{}{\usepackage{subcaption}}
\makeatother

% From https://tex.stackexchange.com/a/645996/211326
%%% apa7 doesn't want to add appendix section titles in the toc
%%% let's make it do it
\makeatletter
\xpatchcmd{\appendix}
  {\par}
  {\addcontentsline{toc}{section}{\@currentlabelname}\par}
  {}{}
\makeatother

%% Disable longtable counter
%% https://tex.stackexchange.com/a/248395/211326

\usepackage{etoolbox}

\makeatletter
\patchcmd{\LT@caption}
  {\bgroup}
  {\bgroup\global\LTpatch@captiontrue}
  {}{}
\patchcmd{\longtable}
  {\par}
  {\par\global\LTpatch@captionfalse}
  {}{}
\apptocmd{\endlongtable}
  {\ifLTpatch@caption\else\addtocounter{table}{-1}\fi}
  {}{}
\newif\ifLTpatch@caption
\makeatother

\begin{document}

\maketitle



\setcounter{secnumdepth}{3}

\setlength\LTleft{0pt}



\ifdefined\Shaded\renewenvironment{Shaded}{\begin{snugshade}\begin{singlespace}}{\end{singlespace}\end{snugshade}}\fi

\setlength{\cslhangindent}{0.5in}

Matriculation Nr.: 401166548

Email:
\href{mailto:Alisoleimani71@gmail.com}{soleimani.ali@stud.hs-fresenius.de}

\section{why to import data?}\label{why-to-import-data}

Data import is the foundational step of any data analysis workflow.
Before cleaning, transformation, modeling, or visualization can occur,
data must be reliably loaded into R. As Wickham and Çetinkaya‑Rundel
(2023) emphasize, \emph{importing data is the gateway to all subsequent
analysis}, and errors at this stage propagate throughout the entire
pipeline.

The R ecosystem provides a rich set of tools for importing data from
diverse formats, including CSV, Excel, JSON, XML, and statistical
software files. The official tidyverse documentation highlights that
choosing the correct import method improves reproducibility,
performance, and data integrity (Posit, 2023).

\section{Why Data import is
Important?}\label{why-data-import-is-important}

Data import is important for three academic reasons:

\textbf{1. Reproducibility}

Reliable import ensures that analyses can be replicated across systems,
operating systems, and environments (Peng, 2011). Incorrect import leads
to inconsistent results.

\textbf{2. Data Integrity}

Parsing errors---such as incorrect column types, encoding issues, or
malformed rows---can introduce bias or invalidate results (Broman \&
Woo, 2018).

\textbf{3. Performance}

Efficient import tools reduce computation time, especially for large
datasets. The fread() function from data.table is shown to outperform
base R by an order of magnitude (data.table Project, 2024).

\section{different data types}\label{different-data-types}

the most common data types encountered in practice include:

1. Text-based tabular data

\begin{itemize}
\tightlist
\item
  CSV\\
\item
  TSV\\
\item
  Fixed-width files
\end{itemize}

2. Spreadsheet data

\begin{itemize}
\tightlist
\item
  Excel (.xls, .xlsx)
\end{itemize}

3. Statistical software formats

\begin{itemize}
\tightlist
\item
  SPSS (.sav)\\
\item
  Stata (.dta)\\
\item
  SAS (.sas7bdat)
\end{itemize}

4. Semi-structured data

\begin{itemize}
\tightlist
\item
  JSON\\
\item
  XML
\end{itemize}

5. Web-based data

\begin{itemize}
\tightlist
\item
  APIs\\
\item
  HTML tables\\
\item
  Remote CSV files
\end{itemize}

6. Databases

\begin{itemize}
\tightlist
\item
  SQL-based systems (via DBI)
\end{itemize}

Each data type requires specialized import tools to preserve structure,
metadata, and encoding.

\section{different packages for each data
type}\label{different-packages-for-each-data-type}

The mapping between data types and R packages is well-established in
academic and official documentation:

{\def\LTcaptype{none} % do not increment counter
\begin{longtable}[]{@{}
  >{\raggedright\arraybackslash}p{(\linewidth - 4\tabcolsep) * \real{0.2568}}
  >{\raggedright\arraybackslash}p{(\linewidth - 4\tabcolsep) * \real{0.3514}}
  >{\raggedright\arraybackslash}p{(\linewidth - 4\tabcolsep) * \real{0.3919}}@{}}
\toprule\noalign{}
\endhead
\bottomrule\noalign{}
\endlastfoot
\textbf{Package} & \textbf{Best For} & \textbf{Why} \\
readr & CSV, TSV, delimited text & Fast, consistent parsing \\
data.table::fread & Large CSV files & Fastest import method \\
readxl & Excel (.xls, .xlsx) & No dependencies, stable \\
haven & SPSS, Stata, SAS & Preserves labels \& metadata \\
jsonlite & JSON, APIs & Handles nested structures \\
xml2 & XML & Robust parsing \\
httr / rvest & Web APIs, HTML scraping & Flexible \\
\end{longtable}
}

\section{Samples of Packages, With Codes for
Each}\label{samples-of-packages-with-codes-for-each}

\subsection{readr (CSV)}\label{readr-csv}

\begin{verbatim}
Rows: 2 Columns: 2
-- Column specification --------------------------------------------------------
Delimiter: ","
dbl (2): x, y

i Use `spec()` to retrieve the full column specification for this data.
i Specify the column types or set `show_col_types = FALSE` to quiet this message.
\end{verbatim}

\begin{verbatim}
# A tibble: 2 x 2
      x     y
  <dbl> <dbl>
1     1     2
2     3     4
\end{verbatim}

\subsection{data.table::fread (Large
CSV)}\label{data.tablefread-large-csv}

\begin{verbatim}
      id value
   <int> <int>
1:     1    10
2:     2    20
3:     3    30
4:     4    40
5:     5    50
\end{verbatim}

\subsection{readxl (Excel)}\label{readxl-excel}

\begin{verbatim}
# A tibble: 3 x 2
  city    population
  <chr>        <dbl>
1 Cologne    1086000
2 Berlin     3769000
3 Hamburg    1841000
\end{verbatim}

\section{haven (SPSS/Stata/SAS)}\label{haven-spssstatasas}

\begin{verbatim}
# A tibble: 3 x 3
     id gender score
  <dbl> <chr>  <dbl>
1     1 male      80
2     2 female    90
3     3 female    85
\end{verbatim}

\subsection{jsonlite (JSON)}\label{jsonlite-json}

\begin{verbatim}
  name age score
1  Ali  25    88
2 Sara  23    91
\end{verbatim}

\subsection{xml2 (XML)}\label{xml2-xml}

\begin{verbatim}
{xml_document}
<employees>
[1] <employee>\n  <id>1</id>\n  <name>Ali</name>\n  <department>Procurement</ ...
[2] <employee>\n  <id>2</id>\n  <name>Sara</name>\n  <department>Finance</dep ...
[3] <employee>\n  <id>3</id>\n  <name>John</name>\n  <department>IT</departme ...
[4] <employee>\n  <id>4</id>\n  <name>Anna</name>\n  <department>HR</departme ...
\end{verbatim}

\section{importing data from multiple
files}\label{importing-data-from-multiple-files}

\section{comparison table based on file
type}\label{comparison-table-based-on-file-type}

{\def\LTcaptype{none} % do not increment counter
\begin{longtable}[]{@{}
  >{\raggedright\arraybackslash}p{(\linewidth - 6\tabcolsep) * \real{0.2500}}
  >{\raggedright\arraybackslash}p{(\linewidth - 6\tabcolsep) * \real{0.2500}}
  >{\raggedright\arraybackslash}p{(\linewidth - 6\tabcolsep) * \real{0.2500}}
  >{\raggedright\arraybackslash}p{(\linewidth - 6\tabcolsep) * \real{0.2500}}@{}}
\toprule\noalign{}
\begin{minipage}[b]{\linewidth}\raggedright
Format
\end{minipage} & \begin{minipage}[b]{\linewidth}\raggedright
Best Package
\end{minipage} & \begin{minipage}[b]{\linewidth}\raggedright
Strengths
\end{minipage} & \begin{minipage}[b]{\linewidth}\raggedright
Limitations
\end{minipage} \\
\midrule\noalign{}
\endhead
\bottomrule\noalign{}
\endlastfoot
CSV & readr & Fast, tidyverse-friendly & Not ideal for huge files \\
Large CSV & fread & Fastest, auto-detects delimiter & Fewer
diagnostics \\
Excel & readxl & No dependencies, reliable & Read-only \\
SPSS/Stata/SAS & haven & Preserves labels & SAS files slower \\
JSON & jsonlite & Handles nested data & Requires JSON knowledge \\
XML & xml2 & Robust parsing & Complex structures \\
Web APIs & httr & Flexible & Requires manual parsing \\
\end{longtable}
}

\section{conclusion}\label{conclusion}

Data import is not merely a technical step but a methodological
foundation for reproducible and reliable analysis. R provides
specialized tools for each data type, and selecting the correct package
ensures optimal performance, data integrity, and workflow efficiency. By
integrating the course notes, R4DS, and CRAN documentation, this handout
provides a comprehensive overview of modern data import practices in R.

\section{Class assignment}\label{class-assignment}

Download the following dataset:
https://raw.githubusercontent.com/allisonhorst/palmerpenguins/main/inst/extdata/penguins.csv

Tasks:

1. Import the dataset using readr, fread, and base R.

2. Measure import time using system.time().

3. Compare column types produced by each method.

4. Write a short paragraph explaining which method is best and why.

\begin{itemize}
\item
  \textbf{Broman, K., \& Woo, K. (2018).} \emph{Data organization in
  spreadsheets.} PLOS Computational Biology. DOI:
  https://doi.org/10.1371/journal.pcbi.1005746
  \href{https://www.bing.com/search?q=\%22https\%3A\%2F\%2Fdoi.org\%2F10.1371\%2Fjournal.pcbi.1005746\%22}{(doi.org
  in Bing)}
\item
  \textbf{data.table Project. (2024).} \emph{data.table: Fast data
  manipulation in R.} CRAN Documentation. (No DOI --- official CRAN
  manual)
\item
  \textbf{Hubchev. (2024).} \emph{Data Science Notes.}
  \url{https://hubchev.github.io/ds/}
\item
  \textbf{Ooms, J. (2014).} \emph{The jsonlite package: A practical and
  consistent mapping between JSON data and R objects.} The R Journal.
  DOI: https://doi.org/10.32614/RJ-2014-011
  \href{https://www.bing.com/search?q=\%22https\%3A\%2F\%2Fdoi.org\%2F10.32614\%2FRJ-2014-011\%22}{(doi.org
  in Bing)}
\item
  \textbf{Peng, R. (2011).} \emph{Reproducible research in computational
  science.} Science. DOI: https://doi.org/10.1126/science.1213847
  \href{https://www.bing.com/search?q=\%22https\%3A\%2F\%2Fdoi.org\%2F10.1126\%2Fscience.1213847\%22}{(doi.org
  in Bing)}
\item
  \textbf{Posit. (2023).} \emph{Importing data with the tidyverse.}
  \url{https://r4ds.hadley.nz/data-import.html} (No DOI --- official
  Posit documentation)
\item
  \textbf{R Project. (2024).} \emph{Read data files from various
  sources.} CRAN Task View. (No DOI --- official CRAN documentation)
\item
  \textbf{RStudio. (2024).} \emph{readxl: Read Excel Files.} CRAN
  Documentation. (No DOI --- official CRAN manual)
\item
  \textbf{Wickham, H. (2014).} \emph{Tidy data.} Journal of Statistical
  Software. DOI: https://doi.org/10.18637/jss.v059.i10
  \href{https://www.bing.com/search?q=\%22https\%3A\%2F\%2Fdoi.org\%2F10.18637\%2Fjss.v059.i10\%22}{(doi.org
  in Bing)}
\item
  \textbf{Wickham, H., \& Çetinkaya‑Rundel, M. (2023).} \emph{R for Data
  Science (2nd ed.).} Posit. \url{https://r4ds.hadley.nz} (No DOI ---
  open textbook)
\end{itemize}

:::






\end{document}
